\chapter{Radiosity}\label{ch:radiosity}

In order to achieve global illumination, the Rendering Equation (Equation~\ref{eq:rendering_equation}) must be solved. Three main approaches exist for solving it:
\begin{itemize}
    \item \ul{Path tracing}: it is a Monte Carlo method that estimates the solution of the rendering equation by randomly sampling light paths. It was explained in Section~\ref{sec:path_tracing}.
    \item \ul{Radiosity}: it will be explained in detail in this chapter.
    \item \ul{Photon mapping}: it will be explained in detail in Chapter~\ref{ch:photon_mapping}.
\end{itemize}

In order to understand the Radiosity method, some radiometric quantities need to be defined first. Radiometry is the science of measuring electromagnetic radiation, including visible light. It provides a set of quantities that are used to describe the properties of light and how it interacts with surfaces.
\begin{definicion}[Radiant Energy]
Radiant energy is the total amount of energy that is carried by a beam of light.
\begin{equation*}
    Q\qquad [J]
\end{equation*}
\end{definicion}

\begin{definicion}[Radiant Flux]
    Also known as Radiant Power, it is the rate at which radiant energy is transferred or emitted. It is defined as the amount of radiant energy per unit time.
\begin{equation*}
    \Phi = \frac{dQ}{dt} \qquad [W=\nicefrac{J}{s}]
\end{equation*}
\end{definicion}

\begin{definicion}[Irradiance]
Irradiance is the amount of radiant flux that is incident on a surface per unit area.
\begin{equation*}
    E = \frac{d\Phi}{dA} \qquad [\nicefrac{W}{m^2}]
\end{equation*}
\end{definicion}

\begin{definicion}[Radiosity]
Radiosity is the amount of radiant flux that is outgoing from a surface per unit area.
\begin{equation*}
    B = \frac{d\Phi}{dA} \qquad [\nicefrac{W}{m^2}]
\end{equation*}
\end{definicion}

\begin{observacion}
    The difference between irradiance and radiosity is that the former measures the incoming light on a surface, while the latter measures the outgoing light from a surface. In other words, irradiance is the amount of light that hits a surface, while radiosity is the amount of light that leaves a surface.
\end{observacion}

For the next measures, the solid angle is needed. It is a measure of the amount of space that an object occupies in the field of view of an observer. It is defined as the area of the projection of the object onto a unit sphere centered at the observer. More generally, if the sphere used for the projection has a radius $r$, and the surface area of the projection is $A$, the solid angle $\Omega$ is defined as:
\begin{equation*}
    \Omega = \frac{A}{r^2} \qquad [sr]
\end{equation*}
The unit of solid angle is the steradian ($sr$) (an adimensional unit), which is defined as the solid angle subtended by an area of $r^2$ on the surface of a sphere of radius $r$. It is the three-dimensional equivalent of the two-dimensional radian.

\begin{definicion}[Radiant Intensity]
Radiant intensity is the amount of radiant flux that is emitted by a source in a particular direction per unit solid angle.
\begin{equation*}
    I = \frac{d\Phi}{d\vec{\omega}} \qquad [\nicefrac{W}{sr}]
\end{equation*}
\end{definicion}

\begin{definicion}[Radiance]
Radiance is the amount of radiant flux that is emitted, reflected, transmitted or received by a surface per \emph{projected} unit area and per unit solid angle.
\begin{equation*}
    L = \dfrac{dI}{dA \cdot \cos(\theta)} =
    \frac{d^2\Phi}{dA \cdot \cos(\theta) \cdot d\vec{\omega}} \qquad [\nicefrac{W}{m^2 \cdot sr}]
\end{equation*}
where $\theta$ is the angle between the normal of the surface and the direction of the incoming light. The term $\cos(\theta)$ is used to account for the fact that the effective area of the surface that is illuminated by the light decreases as the angle of incidence increases.
\end{definicion}
Radiance is the most used quantity in computer graphics, because it is the one that is conserved\footnote{In fact, it only happens in vacuum.} along a ray of light. This means that the radiance of a ray of light does not change as it travels through space, unless it interacts with a surface. This property makes it easier to compute the illumination of a scene, because we can trace rays of light from the camera to the scene and compute the radiance at each point of intersection.

\begin{observacion}
    So many measures may be confusing. The basic one is the radiant energy, and then the radiant flux is the rate of change of the radiant energy. All of the other measures are derived from the flux energy connecting it with other quantities:
\begin{itemize}
    \item Irradiance and Radiosity: Flux per unit area.
    \item Radiant Intensity: Flux per unit solid angle.
    \item Radiance: Flux per projected unit area and per unit solid angle.
\end{itemize}
\end{observacion}

Two essential connections between Radiance and Irradiance and between Radiance and Radiosity are given by the following formulas:
\begin{align}\label{eq:irradiance_radiance}
    E(x) &= \int_{\Omega} L_i(x, \vec{\omega}) \cdot \cos(\theta) \cdot d\vec{\omega} \\
    B(x) &= \int_{\Omega} L_o(x, \vec{\omega}) \cdot d\vec{\omega}
\end{align}
where $\Omega$ is the hemisphere of directions above the point $x$, and $\theta$ is the angle between the normal of the surface at point $x$ and the direction $\vec{\omega}$.

\begin{definicion}[Reflectance]
    Reflectance is the fraction of incident light that is scattered when it hits a surface (that is not absorbed).
    \begin{equation*}
        \rho(x) = \frac{\text{Reflected Flux}}{\text{Incident Flux}} \qquad [\text{adimensional}]
    \end{equation*}
\end{definicion}

This last material property connects the irradiance and the radiosity of a surface with the following formula:
\begin{equation}\label{eq:radiosity_irradiance}
    B(x) = \rho(x) \cdot E(x)
\end{equation}

Once all these concepts have been defined, the Radiosity method can be explained.

\section{The Radiosity Equation}

In order to solve the rendering equation using the Radiosity method, it needs to be translated into radiosity terms. Three first changes are needed:
\begin{itemize}
    \item \ul{Iterating over points instead of directions}:
    
    Instead of iterating over all the possible $\vec{\omega}'$ directions of the $\Omega$ hemisphere above the surface, it iterates over all the points $y$ in the scene $S$. These points are infinitelly small with area $dA_y$.
    \item \ul{Visibility factor}:
    
    Since it iterates over all the points $y$ in the scene, some of them will not be visible from point $x$ (that is, there will be an object between them). Therefore, only the points $y$ that are visible from $x$ should be considered. In order to do that, a visibility factor $V(x, y)$ is included, which is 1 if the points $x$ and $y$ are visible to each other (that is, there is no object between them), and 0 otherwise.
    \begin{equation*}
    V(x, y) = \begin{cases}
        1 & \text{if $x$ and $y$ are visible to each other} \\
        0 & \text{otherwise}
    \end{cases}
    \end{equation*}
    \begin{observacion}
        Since we consider points $y$ with an infinitesimal area, the visibility factor is either 0 or 1. However, if we consider patches with a finite area, the visibility factor can take any value between 0 and 1, depending on the fraction of the patch that is visible from point $x$.
    \end{observacion}

    \item \ul{Geometry factor}:
    
    A change of variable from ${\omega}'$ to $A_y$ is needed. From the definition of solid angle, we have that:
    \begin{equation*}
        \omega' = \frac{A_{y,\text{projected}}}{r^2}
        \Longrightarrow
        d\omega' = \frac{dA_{y,\text{projected}}}{r^2}
    \end{equation*}
    where $r=\|x-y\|$ is the distance between points $x$ and $y$, and $A_{y,\text{projected}}$ is the area of the projection of the patch at point $y$ onto a plane tangent to the sphere of radius $r$ centered at point $x$ used for the solid angle. That projection is given by:
    \begin{equation*}
        A_{y,\text{projected}} = A_y \cdot \cos(\theta_y)
    \end{equation*}
    where $\theta_y$ is the angle between the normal of the surface at point $y$ ($N'$) and the vector from $y$ to $x$, which is $-\omega'$. Therefore, we have that:
    \begin{equation*}
        A_{y,\text{projected}} = A_y \cdot \cos(\theta_y) = A_y \cdot (N' \cdot (-\omega'))
    \end{equation*}

    Taking this into account, we have:
    \begin{equation*}
        d\omega' = \frac{dA_{y,\text{projected}}}{r^2} = \frac{dA_y \cdot (N' \cdot (-\omega'))}{\|y-x\|^2}
    \end{equation*}

    In order to have a simpler formula, a geometry factor $G(x, y)$ is defined so that:
    \begin{equation*}
    G(x, y) dA_y = (\omega' \cdot N) \ d\omega'
    \Longrightarrow
    G(x, y) = \frac{(\omega' \cdot N) \cdot (N' \cdot (-\omega'))}{\|y-x\|^2}
    \end{equation*}
    
    This geometry factor accounts for the orientation of the surfaces at points $x$ and $y$, as well as the distance between them.
\end{itemize}

Taking these three changes into account, the rendering equation can be rewritten\footnote{It should be noted that it is simply rewritten, not approximated.} as:
\begin{align*}
    L(x, \lambda, \omega) &= L_e(x, \lambda, \omega) + \int_{y\in S} f_r(x, \lambda, \omega', \omega) \cdot L_i(y, \lambda, \omega') \cdot V(x, y) \cdot G(x, y) \ dA_y
\end{align*}

\subsection{Diffuse Reflection Simplification}

The radiosity method is based on the assumption that all surfaces in the scene are perfectly diffuse reflectors, which means that they reflect light equally in all directions. This simplification allows to further simplify the rendering equation, as the BRDF of a perfectly diffuse reflector is constant and equal to:
\begin{equation*}
    f_r(x, \lambda, \omega', \omega) = \frac{\rho(x)}{\pi}
\end{equation*}
% // TODO: Añadir demostración
where $\rho(x)$ is the reflectance of the surface at point $x$. Therefore, the rendering equation can be rewritten as:
\begin{align*}
    L(x, \lambda, \omega) &= L_e(x, \lambda, \omega) + \int_{y\in S} \frac{\rho(x)}{\pi} \cdot L_i(y, \lambda, \omega') \cdot V(x, y) \cdot G(x, y) \ dA_y
\end{align*}

In order to simplify the notation even more, a \emph{transfer function} $G'(x, y)$ is defined as:
\begin{equation*}
    G'(x, y) = \frac{V(x, y) \cdot G(x, y)}{\pi}
\end{equation*}
so the rendering equation can be rewritten as:
\begin{align*}
    L(x, \lambda, \omega) &= L_e(x, \lambda, \omega) + \rho(x) \cdot \int_{y\in S} L_i(y, \lambda, \omega') \cdot G'(x, y) \ dA_y
\end{align*}

\subsection{Translating Radiance to Radiosity}

Changing the perspective from radiance to radiosity, the resulting equation is:
\begin{align*}
    B(x) &= B_e(x) + \rho(x) \cdot \int_{y\in S} B(y) \cdot G'(x, y) \ dA_y
\end{align*}
where $B_e(x)$ is the emitted radiosity at point $x$.

\subsection{The Discretized Radiosity Equation}

In order to solve the radiosity equation, the scene is divided into small patches with:
\begin{itemize}
    \item Finite area $A_i$.
    \item Constant radiosity $B_i$.
    \item Constant reflectance $\rho_i$.
    \item Constant emitted radiosity $B_i^e$.
\end{itemize}

In addition, the integral of the geometry factor is approximated by a form factor $F_{ij}$, which is defined as:
\begin{equation*}
    F_{ij} = \frac{1}{A_i} \int_{x\in A_i} \int_{y\in A_j} G'(x, y) \ dA_y \ dA_x
\end{equation*}
This form factor represents the fraction of the radiosity that leaves patch $i$ and arrives at patch $j$. Therefore, the radiosity equation can be rewritten as:
\begin{equation}\label{eq:radiosity_equation}
    B_i = B_i^e + \rho_i \cdot \sum_{j=1}^N B_j\cdot F_{ij}
\end{equation}
where $N$ is the total number of patches in the scene.

\subsubsection{The ``direction'' of the Form Factors}

The reader could think that, if $F_{ij}$ represents the fraction of the radiosity that leaves patch $i$ and arrives at patch $j$, then the Equation~\ref{eq:radiosity_equation} should use $F_{ji}$ instead of $F_{ij}$. In order to understand why $F_{ij}$ is used, the following connection between both form factors is needed:
\begin{equation*}
    A_i \cdot F_{ij} = A_j \cdot F_{ji}
\end{equation*}

From the Equation~\ref{eq:radiosity_equation}, if every term is multiplied by $A_i$ (representing the total radiant flux leaving patch $i$), we have:
\begin{equation*}
    A_i \cdot B_i = A_i \cdot B_i^e + \rho_i \cdot \sum_{j=1}^N B_j\cdot A_i \cdot F_{ij}
\end{equation*}

If the connection between the form factors is applied, we have:
\begin{equation*}
    A_i \cdot B_i = A_i \cdot B_i^e + \rho_i \cdot \sum_{j=1}^N B_j\cdot A_j \cdot F_{ji}
\end{equation*}
which is the reasonable interpretation of the equation: the total radiant flux leaving patch $i$ is equal to the emitted radiant flux from patch $i$ plus the reflected radiant flux from patch $i$, which is the portion of the radiant flux that arrives at patch $i$ from all the other patches $j$ (which is $B_j\cdot A_j\cdot F_{ji}$) multiplied by the reflectance $\rho_i$. Therefore, when working with radiant flux, the form factor $F_{ji}$ may be used having a ``reasonable'' interpretation of the equation, but when working with radiosity, the form factor $F_{ij}$ is needed.

\subsubsection{Obtaining the Form Factors}

The definition of the form factors was given by:
\begin{equation*}
    F_{ij} = \frac{1}{A_i} \int_{x\in A_i} \int_{y\in A_j} G'(x, y) \ dA_y \ dA_x
\end{equation*}

This is the most computationally expensive part of the radiosity method, because $G'(x, y)$ is a complex function that depends on the geometry of the scene and the visibility between patches $i$ and $j$. Therefore, several methods have been proposed to approximate the form factors.
\begin{itemize}
    \item \ul{Hemicube method}: it is a method that aproximates the foam factor between a differential patch $dA_i$ and a patch $A_j$ by projecting the patch $A_j$ onto a hemispherical sphere centered at the point $x$ of the patch $dA_i$.
    \begin{itemize}
        \item The projection onto the sphere takes into account $N' \cdot (-\omega')$ and $\|y-x\|^2$.
        \item The projection onto the circle takes into account $N \cdot \omega'$ and $\pi$.
    \end{itemize}

    However, projecting the patch $A_j$ onto a hemispherical sphere is computationally expensive. Therefore, the hemicube method approximates the projection onto the hemisphere by projecting the patch $A_j$ onto a hemicube centered at the point $x$ of the patch $dA_i$. The projection onto the hemicube is much faster than the projection onto the hemisphere, as it can be implemented rasterizing from $x$ using 5 different viewing directions (one for each face of the hemicube).
\end{itemize}

\section{The Radiosity Method}

Once the radiosity equation is obtained, let's explain how everythig is combined in order to obtain the final image.
\begin{enumerate}
    \item The scene is divided into small patches with finite properties and constant properties.
    \item The form factors between all the patches are calculated using one of the different methods that exist for that purpose, as the hemicube method. It is the most computationally expensive part of the radiosity method, as obtaining the visibility between patches and the geometry factor is complex.
    
    It should be noted that the form factors only depend on the geometry of the scene and the visibility between patches, but not on the material properties of the patches. Therefore, if the geometry of the scene does not change, the form factors can be precomputed and reused for different material properties.
    \item For each patch and each component of the light\footnote{In practice, the light is usually divided into three components: red, green and blue. However, theoretically, it could be computed for each wavelength of light.}, the radiosity for that patch and that component of the light is calculated. Writing the Radiosity Equation (Equation~\ref{eq:radiosity_equation}) in matrix form, we have:
    \begin{equation*}
        \begin{pmatrix}
            B_1 \\
            B_2 \\
            \vdots \\
            B_N
        \end{pmatrix}
        = \begin{pmatrix}
            B_1^e \\
            B_2^e \\
            \vdots \\
            B_N^e
        \end{pmatrix}
        + \begin{pmatrix}
            \rho_1 0 & \cdots & 0 \\
            0 & \rho_2 & \cdots & 0 \\
            \vdots & \vdots & \ddots & \vdots \\
            0 & 0 & \cdots & \rho_N
        \end{pmatrix}
        \begin{pmatrix}
            F_{11} & F_{12} & \cdots & F_{1N} \\
            F_{21} & F_{22} & \cdots & F_{2N} \\
            \vdots & \vdots & \ddots & \vdots \\
            F_{N1} & F_{N2} & \cdots & F_{NN}
        \end{pmatrix}
        \begin{pmatrix}
            B_1 \\
            B_2 \\
            \vdots \\
            B_N
        \end{pmatrix}
    \end{equation*}

    Defining a matrix $T$ called the transfer matrix as:
    \begin{equation*}
        T = \begin{pmatrix}
            \rho_1 0 & \cdots & 0 \\
            0 & \rho_2 & \cdots & 0 \\
            \vdots & \vdots & \ddots & \vdots \\
            0 & 0 & \cdots & \rho_N
        \end{pmatrix}
        \begin{pmatrix}
            F_{11} & F_{12} & \cdots & F_{1N} \\
            F_{21} & F_{22} & \cdots & F_{2N} \\
            \vdots & \vdots & \ddots & \vdots \\
            F_{N1} & F_{N2} & \cdots & F_{NN}
        \end{pmatrix}
    \end{equation*}
    the radiosity equation can be rewritten as:
    \begin{equation*}
        B = B^e + T \cdot B
    \end{equation*}

    There are two main methods for solving this equation and obtaining the radiosity of each patch:
    \begin{itemize}
        \item \ul{Algebraic method}: it consists in rearranging the equation as:
        \begin{align*}
            B &= B^e + T \cdot B \\
            (I - T) \cdot B &= B^e \\
            B &= (I - T)^{-1} \cdot B^e
        \end{align*}
        where $I$ is the identity matrix. However, this method is computationally expensive, as it requires inverting the matrix $I - T$, which has a complexity of $O(N^3)$, where $N$ is the number of patches in the scene.
        \item \ul{Iterative method}: it consists in iteratively applying the equation until convergence. The equation can be solved iteratively as:
        \begin{align*}
            B^{(0)} &= B^e \\
            B^{(k+1)} &= E + T \cdot B^{(k)}
        \end{align*}
        where $B^{(k)}$ is the radiosity at iteration $k$, representing the light until $k$ bounces.
        
        This method is more efficient than the algebraic method, as it has a complexity of $O(N^2)$ per iteration, and it usually converges in a few iterations.
    \end{itemize}
    \item Once the radiosity of each patch is obtained, a radiosity value needs to be assigned to each vertex of the scene, as the final image is rendered using rasterization. There are two main methods for doing that:
    \begin{itemize}
        \item \ul{Flat Method}: it consists in assigning the radiosity of the patch to all the vertices of the patch. This method is simple, but it can produce artifacts in the edge of the patches, as the radiosity changes abruptly between patches.
        \item \ul{Smooth Method}: it consists in assigning the radiosity of each patch to the vertices of the patch, and then averaging the radiosity of the vertices that belong to multiple patches. This method produces smoother results, as it avoids abrupt changes in the radiosity between patches.
        
        However, it is more computationally expensive than the flat method, and it needs to take into account the normals, as there are cases where real edges exist in the scene.
    \end{itemize}
    \item Finally, the final color for the corresponding light component of each vertex is calculated from the radiosity as follows:
    \begin{equation*}
        C = \frac{B}{B+1}
    \end{equation*}
    This formula is used to convert the radiosity, which can take any value in $[0, +\infty[$, to a color value in $[0, 1[$. The formula is designed so that when the radiosity is low, the color is close to 0, and when the radiosity is high, the color approaches 1.
    \item Once the color of each vertex is obtained, the final image can be rendered using rasterization.
\end{enumerate}